\section{流程清单与验证}
\subsection{训练前准备}
在进入 self-rewarding 训练前，Weston 的建议包括：选择 Self-Instruct seeds、准备 SelfT evaluators、配置 inference compute budget、设定 cross-check 频率。这些步骤让每一次迭代都能受控地生成 verified response，其中 ``Self-Instruct 产生的任务必须覆盖数学、编程、判断''（SRT 674-676）。

\begin{itemize}
    \item 设定 Self-Instruct seed（Math / Code / Reasoning / Dialogue）并扩写成 diverse prompts。
    \item 准备 verified response pipeline，制定 cross-check 标签标准。
    \item 预先训练或 fine-tune SelfT evaluator，确保 judge 可区分 draft 与 verified。
    \item 规划 inference-time compute，避免 cross-check 输出成为 compute bottleneck（SRT 3236-3252）。
    \item 备份 Self-Feeding 对话数据，用于后续监督。
\end{itemize}

\subsection{评估后处理}
每轮训练结束后，需要把 results 用于 meta judge 校准、统计 win rate，并记录 verified response 是否跨 iteration 一致。Weston 特别提到 ``kept verifying itself and iterating'' 的过程（SRT 1908-2027），建议把 verified response 抽出与 Alpaca Eval 2 比对，并用 judge 重新评分，确保 plateau 被打破。

\begin{itemize}
    \item 让 meta judge 对 draft 与 verified response 再次打分，实时捕捉 drift。
    \item 把最新 verified response 与 Alpaca Eval 2 / Alpaca Eval 2 baseline 比较。
    \item 汇总 error logs，特别关注 hallucination 与 overthinking 的样本。
    \item 依据 judge 反馈调整 reward model 的 loss 权重与 cross-check 频率。
    \item 将高质量 verified response 标记为 future seed，持续扩充 Self-Instruct 语料。
\end{itemize}

\subsection{本章小结}
训练前的准备负责保证不同任务与 judge 都在位，评估后处理则负责纠偏与总结：只有把 verified response 既拿去训练又拿去评估，才能让 self-rewarding 持续推动 reasoning 质量。

